\section{Consequences, distinctions, and proof boundaries}\label{sec:scope}
The general conclusion has three operationally different ingredients: the full-history acquisition risk $B_N$, the acquired geometric allocation $\Psi_N$, and the forced transport operator in \eqref{eq:transport}. They originate in the same raw experiment but are not interchangeable. A formal parameter dimension controls none of them by itself. In the singular sensor, for example, the raw target may carry coordinates killed by the report map, the absorbed law may mix dimensions, and its total mass is only $1-(1-s)^N$.

The two realizations verify different causal mechanisms. The filter has positive continuous likelihoods, repeated Bayesian normalization, and contraction in log-odds; its acquired density is established by an actual report Jacobian. The singular sensor has an absorbing reveal, exact holds of arbitrarily long duration, and an acquired measure with intersections and atoms; its transport is a disjoint indicator identity. Neither calculation is a premise of Theorem~\ref{thm:main}. The main proof consists of the predictive realization, actual-mass lower estimate, global valid cover, and acquisition-weighted transport argument, followed by common-task morphism composition.

\subsection{Relation to existing theory}
The checkpoint identity is the conditional-mean projection underlying quantization, not a new alternative to it. The anisotropic box proof is included to expose its constants through vanishing physical widths and to accommodate actual component weights and discarded components. The general study of quantization of singular probabilities is much broader than this box lemma \cite{graf}; local error estimates for Ahlfors--David measures provide another developed direction \cite{zhu}. We do not claim that the present paper supersedes those results.

The difference from numerical filter quantization \cite{pages,sellami,feuer} is the simultaneous lower bound for the same acquired task and the explicit separation of retained labels, workspace, program data, clock, and raw acquisition. The innovation argument does not charge another rounding on an exactly represented hold. The difference from unforced random-iteration convergence \cite{diaconis} is the weighted second-moment estimate for the forcing introduced by a finite machine. The experiment-comparison argument \cite{blackwell} remains a common-loss argument; subtracting two different oracles would destroy it.

\subsection{Necessary distinctions illustrated by exact examples}
First, the support $[0,1]$ equipped with a point mass and the same envelope equipped with a uniform law have respectively zero one-point distortion and strictly positive finite-label distortion. Hence a reachability envelope is not an acquired law. More generally a component of probability $w$ contributes with weight $w$, even if its geometric diameter is large.

Second, an identity update has no need for repeated quantization. If the initial state has already been replaced by a representative, the same representative can be retained indefinitely. The bound $Nr$ obtained by blindly injecting a covering radius at every hold is not an inherent resource requirement. The acquisition-weighted certificate corrects precisely this failure of accounting.

Third, a statistical garbling can increase memory needed to match its \emph{own} oracle. Observing a fair hidden bit exactly produces only posterior probabilities zero and one. Adding independent Gaussian noise to this observation produces a continuously varying posterior probability. Two labels reproduce the first oracle exactly, but no finite number reproduces the latter continuously distributed oracle exactly under square score. This is compatible with data processing because the baselines differ.

Fourth, let $P_N=\delta_0$ and $Q_N=(1-e^{-Nc})\delta_0+e^{-Nc}\delta_1$, $c>0$. Their total variation distance tends to zero, but the exponential rates at the point one are infinite and $c$, respectively. The common-risk defect of this paper is therefore not a certificate of large-deviation equivalence.

Finally, a hidden mode or an internal simulator register is not free because its name is omitted from a state vector. Independent pairs of retained states have the product cardinality in Theorem~\ref{thm:morphism}. Conversely that product is a constructive upper bound, not a theorem that every simulator requires it. An exact label machine and a finite-bit numerical implementation are also different objects: the latter retains the nonzero calibration, workspace, and readout terms of \eqref{eq:joint}. Exact-arithmetic zero distortion is not a promise of a finite-precision physical output with zero error.

\subsection{Scope of the mathematical conclusion}
The main theorem is a conditional but quantitative class theorem: its certificates do not contain optimal risks, and they are completely checked for the two raw-kernel classes above. It establishes finite-horizon Bayesian prediction curves under fixed exploration, and common-task causal risk transfer over whatever strategy class has a uniform simulator certificate. It does not prove an unrestricted optimal exploration theorem, a parameter-uniform predictive quotient for every prior, or the learning of an unknown kernel from an unpaid pilot. The intersecting-strata result is uniform in distances and widths for finitely many uniformly bi-Lipschitz acquired charts, at the stated budget threshold. General singular laws without this geometry require new actual-mass estimates, not a formal change of notation.

The remaining extensions are specific mathematical questions. One is to derive acquisition-weighted moment bounds for recurrent nonuniform filters that may expand during active updates, beyond the exact-hold and contractive-active mechanisms proved here. Another is to derive uniform charts and optimal calibration acquisition costs from unknown kernel classes. A third is to match the complete resource vector, including growing controller and simulator states, without assuming constant overhead. Infinite-horizon path-law equivalence, large deviations, and unbounded-generator closure require their own stronger certificates. None is inferred from the finite-time theorem or delegated to a realization as an unproved foundation.
